\section{The extension and its exact scope}\label{sec:scope}

The current repository is already far beyond the proposal that continuous
objects should be replaced by finite certificates. Its November 20 packets
include quantitative completion, finite nets of actual Lipschitz functions,
continuous model semantics, holomorphic gluing, coercive operators, and
topological descent. Its \texttt{Derivations} tree contains substantial
\emph{object-theory} arithmetic, including rational sums, weighted error
budgets, and finite conditional-expectation identities
\cite{repo-overview,repo-derivations,repo-abstract,repo-continuous-two}.
The overview's older packet counts should therefore not be treated as a
complete inventory. We inspected the pinned source tree directly.

This report extends that work in a direction where the previous abstractions
need additional structure: a small error in mean, a small uniform error, and
a small operator error support different mathematical operations. Treating
them interchangeably can turn a correct finite computation into an invalid
continuous conclusion. Conversely, demanding uniform error everywhere can
make useful constructions needlessly expensive or impossible.

\subsection{What is reused and what is added}

\begin{longtable}{@{}p{0.23\textwidth}p{0.34\textwidth}p{0.36\textwidth}@{}}
\toprule
Area & Existing material inspected & Contribution here\\
\midrule
\endhead
Arithmetic and proof composition & \texttt{PRA.qfInd}, rational arithmetic and finite sums, finite Fubini, proof-carrying error DAGs & A concrete route for clipping, mean-error, and kernel estimates through those interfaces; the arithmetic infrastructure is credited rather than reintroduced as new.\\
Completion and topology & Regular names, dense presentations, compact nets, coherent witnesses & Norm-sensitive extension: $L^1$, convergence in measure, strong operator action, and exact chain identities have distinct interfaces.\\
Expectation & Finite natural weights, common denominators, conditional cell regrouping, a dominated truncation bound & Stable clipping-Vitali estimate, an $L^1$ limit constructor from Cauchy-in-measure/UI data, and completion of finite conditional averages.\\
Category theory & Presented setoids, completion functors, finite diagrams and gluing & A composable category of probability-preserving operators; its tensor and identity are justified in the appropriate norms.\\
Model theory & Quantitative transfer for compact continuous formulas with infimum and supremum & Integration binders and mixed error budgets; preservation of binder order and explicit limits of $L^1$ transfer.\\
Complex analysis & Certified power-series and derivative tails & Integration of $L^1$-valued series, differentiation under the integral, and a complete rational example.\\
Algebraic topology & Finite nerve and chain comparisons; a separate $\mathbb F_2$ perturbation calculation & Signed normed perturbation with a finite nilpotence clock or a Neumann gap, quantitative residuals, and stability.\\
\bottomrule
\end{longtable}

The theory of Markov categories, probability algebras, quantitative convergence,
and the homological perturbation lemma are established subjects
\cite{fritz,berenstein-henson,neri-pischke,avigad-dean-rute,crainic2004}.
The contribution is their explicit, compatible use within this repository's
first-order presentation discipline, together with the correction and
additional quantitative constructions below. We do not claim a new classical
Vitali theorem or a new homological perturbation lemma.

\subsection{A specific mathematical repair}

The Vitali proof in
\texttt{expected\_value\_ipra\_2026-11-16}, Section 12, uses the pointwise claim
\begin{equation}\label{eq:false-old-vitali}
 |X-Y|\le |X-Y|\mathbf1_{|X|\le K,\,|Y|\le K}
       +|X|\mathbf1_{|X|>K}+|Y|\mathbf1_{|Y|>K}.
\end{equation}
On the one-point probability space, set $K=1$, $X=2$, $Y=-1$.
The two sides are $3$ and $2$, respectively. This is a defect in the
displayed proof, not a counterexample to the stated Vitali conclusion
\cite{repo-expectation-packet}. The separate checked
\texttt{Probability.TruncationEstimate} uses a common dominating function
and a different, valid three-region estimate; it does not assert
\eqref{eq:false-old-vitali} \cite{repo-truncation}.

The replacement uses
\[
 c_K(t)=\max(-K,\min(t,K)),\quad
 t_K(X)=\int (|X|-K)_+,\quad
 \rho(X,Y)=\int\min(1,|X-Y|).
\]
For $K\ge1$ we prove
\begin{equation}\label{eq:main-vitali-intro}
 \norm{X-Y}_1\le t_K(X)+t_K(Y)+2K\rho(X,Y).
\end{equation}
All three quantities on the right are stable under $L^1$ approximation.
Unlike a hard threshold probability, they require no decision about equality
at a boundary. The theorem subsequently constructs the limit from pairwise
estimates and explicit moduli; it does not presuppose that the desired limit
has already been represented in $L^1$.

\subsection{How to read the claims}

There are three distinct deliverables. First, this report proves statements
about supplied presentations and quantitative hypotheses. Second, the Lean
project proves selected finite arithmetic and logical kernels. Third, the
object-language plan specifies the numerical functions, quantifier-free
invariants, existing interfaces, and missing proof values needed to export
the new results to \texttt{Prf PRAAx}. A checked host theorem does not itself
complete that third task.

The infinite notation is a readable account of numerical constructions.
The iPRA content is always a theorem about finite observations at arbitrary
natural-number precisions. Each evaluator, refinement operation, tail
modulus, and clock used in a uniform claim must be supplied as a fixed PR
definition, or be implemented by a bounded interpreter with a separately
proved clock. Neither arbitrary real functions nor arbitrary measures are
being quantified over as objects of first-order iPRA.

The ZIP includes editable sources, a compiled PDF, the complete standalone
Lean sources, exact-rational reference programs, an immutable source map,
and a claim-status ledger. No remote repository files were changed.
